\section{Pseudorandom functions}\label{sec:prelims-prfs}

Pseudorandom functions enable a party holding a secret key to evaluate a
function whose outputs are indistinguishable from those of a truly random
function to anyone without the key.

\paragraph{Syntax.} A pseudorandom function family consists of algorithms
defined as follows:
%
\begin{description}

  \item[$(\pp, \key) \sample \PRFgen(\secparam)$ :] The key generation algorithm
takes the security parameter as input, and outputs public parameters $\pp$
together with a key $\key$.

  \item[$y \leftarrow \PRFeval(\pp, \key, x)$ :] The evaluation algorithm takes
public parameters $\pp$, key $\key$, and input $x$ in the domain, and outputs
$y$ in the range.

\end{description}
%
We say $\PRF = (\PRFgen, \PRFeval)$ is a pseudorandom function family with
domain $\domain$ and range $\range$ if $\PRFeval(\pp, \key, \cdot)$ is
efficiently computable for every $(\pp, \key)$ output by $\PRFgen(\secparam)$,
and if it satisfies pseudorandomness, which we define below. We use the standard
syntax and security notions for pseudorandom functions~\cite{KatLin14}.

The public parameters are carried explicitly because the algebraic
instantiations of later chapters need them: the domain, the range, or the
description of the function itself may depend on group elements that are common
to every key, and those elements are public. The textbook setting is the special
case $\pp = \varnothing$, and we suppress $\pp$ where it plays no role.

\paragraph{Pseudorandomness.} Pseudorandomness requires that no efficient
adversary, given the public parameters and oracle access to the function, can
distinguish it from a function sampled uniformly at random from all functions
$\domain \to \range$.

\begin{experimentbox}{Pseudorandomness}{prf}

  \begin{enumerate}

    \item The challenger samples $b \sample \{0,1\}$ and generates $(\pp, \key)
\sample \PRFgen(\secparam)$. If $b = 1$, it additionally samples a function $R$
uniformly at random from all functions $\domain \to \range$. It sends $\pp$ to
the adversary.

    \item The adversary is given oracle access to $\mathcal{O}_{\PRF}$, which on
input $x \in \domain$ returns $\PRFeval(\pp, \key, x)$ if $b = 0$, and $R(x)$ if
$b = 1$.

    \item The adversary outputs a bit $b'$.

  \end{enumerate}

  $\adv$ wins the experiment if $b' = b$.

\end{experimentbox}

The pseudorandomness advantage of $\adv$ is defined as follows:
%
\begin{equation*}
%
  \Adv_{\PRF,\adv}^{\mathsf{prf}}(\secparam) = \left| \Pr[b'=b]-\frac{1}{2}
\right|.
%
\end{equation*}
%
We say $\PRF$ is a secure pseudorandom function if this advantage is small for
every admissible adversary $\adv$ making a polynomial number of queries to
$\mathcal{O}_{\PRF}$.

\paragraph{Public parameters and pseudorandomness.} Because $\pp$ is public, it
is handed to the adversary in the experiment above, and pseudorandomness is
required to hold in its presence. This is not a formality where the parameters
are derived from long-lived secrets: a construction in which some public value
is a deterministic function of material related to the key must still be
pseudorandom to a distinguisher who holds that value. The tag function
of~\Cref{chap:txid} is of exactly this shape. Its public parameter is the user's
certified public tracing key, fixed for the lifetime of the user, and the
security statements of that chapter are made against a distinguisher given the
whole of it.

The same point constrains what may be placed in the \emph{output} of a
pseudorandom function. A function whose output has a component that is identical
on every input is not pseudorandom, whatever the hardness assumption behind it,
since a random function repeats nothing. Material that a later computation needs
but that does not vary with the input therefore belongs in the public
parameters, not in the output. \Cref{sec:txid-construction} makes this
arrangement explicit for the pair of functions used there.
